% ============================================================
% The Autocorrelation Formula
% ============================================================

\documentclass[12pt]{amsart}

\usepackage{amsmath,amssymb,amsthm}
\usepackage{booktabs}
\usepackage{microtype}
\usepackage{url}

\newtheorem{theorem}{Theorem}
\newtheorem{lemma}[theorem]{Lemma}
\newtheorem{proposition}[theorem]{Proposition}
\newtheorem{corollary}[theorem]{Corollary}
\theoremstyle{definition}
\newtheorem{definition}[theorem]{Definition}
\theoremstyle{remark}
\newtheorem{remark}[theorem]{Remark}

\title{The Autocorrelation Formula}

\author{Alexander S.\ Petty}

\email{alexander.petty@gmail.com}
\date{January 2022 (revised April 2026)}
\subjclass[2020]{11A63, 11B83, 42A16}

\begin{document}
\begin{abstract}
We derive an exact character-sum formula for the cyclic
autocorrelation $R(\ell)$ of the base repetend of $1/p$. For
primes where $b$ is a primitive root, the formula is
\[
R(\ell) = \frac{1}{p}
\sum_{k=0}^{p-1}
G(k,\, -k b^{-\ell} \bmod p),
\]
where $G(k,k') = \sum_{d} \widehat{\mathbf{1}}_{B_d}(k)\,
\widehat{\mathbf{1}}_{B_d}(k')$ is the cross-spectral function
of the bin partition. For primitive-root primes, this closes the formula chain
from the digit function to the eigenvalue spectrum: the
bin sizes determine $G$, $G$ determines $R$, and $R$
determines the eigenvalues via the DFT.

The derivation uses the orthogonality of additive characters on
$\mathbb{Z}/p\mathbb{Z}$~\cite{iwaniec,hardy} and the identity
$\sum_{k,k'} G(k,k') = 0$, which forces the off-diagonal
correction to cancel exactly. The cross-spectral function $G$
extends the spectral power $\Phi(k) = G(k,-k)$
from~\cite{paper8} to a two-variable object encoding both
magnitude and phase information of the bin Fourier transforms.
\end{abstract}
\maketitle

% ============================================================
\section{Introduction}

The cross-alignment matrix $\mathbf{A}(p)$ of a prime $p$
in base $b$ encodes the digit-equality structure of the
fractional field~\cite{paper6}. Its eigenvalue spectrum,
which records the harmonic content of the field, was
identified in~\cite{paper7}: for primes where $b$ is a
primitive root modulo $p$, the eigenvalues are the
discrete Fourier transform of the cyclic autocorrelation
$R(\ell)$ of the base repetend $1/p$. The autocorrelation
counts position-wise digit matches between the repetend
and its cyclic rotations.

The spectral power paper~\cite{paper8} computed
$\Phi(k) = \sum_d |\widehat{\mathbf{1}}_{B_d}(k)|^2$ in
closed form as a sum of squared Dirichlet kernels weighted
by the bin sizes. This gave a complete description of the
magnitude data of the bin Fourier coefficients but did not
determine the autocorrelation, since $R(\ell)$ depends on
the phase information that the squaring operation discards.
A gap remained: expressing $R(\ell)$ itself through the
Fourier data of the bin partition.

This paper closes the gap. We introduce the
\emph{cross-spectral function} $G(k, k') = \sum_d
\widehat{\mathbf{1}}_{B_d}(k)\, \widehat{\mathbf{1}}_{B_d}(k')$
of the bin partition, a two-variable extension of $\Phi$
whose antidiagonal $G(k, -k)$ recovers $\Phi(k)$ and whose
off-diagonal entries carry the phase information needed
for $R(\ell)$. We prove (Theorem~\ref{thm:autocorrelation})
that for primes where $b$ is a primitive root,
\[
R(\ell) = \frac{1}{p} \sum_{k=0}^{p-1} G\!\left(k,\,
-k\, b^{-\ell} \bmod p\right),
\]
expressing the autocorrelation as a character sum over
$G$ along an orbit path determined by the shift $\ell$.
The proof rests on a single algebraic identity:
$\sum_{k, k'} G(k, k') = 0$ (Lemma~\ref{lem:G-sum-zero}),
which forces a cancellation in the naive expansion that
collapses two terms into one.

Combined with~\cite{paper3,paper7,paper8}, the autocorrelation
formula closes the chain from the bin sizes of the digit
function to the eigenvalue spectrum of the cross-alignment
matrix in finitely many explicit steps. Every step is
computable from $p$ and $b$ alone, with no empirical
counting required. The multi-orbit case (where $b$ is not
a primitive root) is not handled by the present technique
and remains open; we discuss the obstruction in
\S\ref{sec:multi-orbit}.

% ============================================================
\section{The Cross-Spectral Function}

\begin{definition}
The \emph{cross-spectral function} of the bin partition is
\[
G(k,k') = \sum_{d=0}^{b-1}
\widehat{\mathbf{1}}_{B_d}(k)\,
\widehat{\mathbf{1}}_{B_d}(k'),
\]
where $\widehat{\mathbf{1}}_{B_d}(k) = \sum_{r \in B_d}
e^{-2\pi i kr/p}$ is the additive Fourier transform of the $d$-th
bin indicator.
\end{definition}

The antidiagonal $G(k,-k) = \sum_d |\widehat{\mathbf{1}}_{B_d}(k)|^2
= \Phi(k)$ recovers the spectral power from~\cite{paper8}
(using $\widehat{\mathbf{1}}_{B_d}(-k) = \overline{\widehat{\mathbf{1}}_{B_d}(k)}$
for real-valued indicators). The
off-diagonal entries encode phase relationships between bins at
different frequencies.

\begin{lemma}\label{lem:G-sum-zero}
$\sum_{k=0}^{p-1} \sum_{k'=0}^{p-1} G(k,k') = 0$.
\end{lemma}

\begin{proof}
By definition,
\[
\sum_{k, k'} G(k, k')
= \sum_{k, k'} \sum_d
\widehat{\mathbf{1}}_{B_d}(k)\,
\widehat{\mathbf{1}}_{B_d}(k')
= \sum_d \biggl(\sum_{k=0}^{p-1}
\widehat{\mathbf{1}}_{B_d}(k)\biggr)
\biggl(\sum_{k'=0}^{p-1}
\widehat{\mathbf{1}}_{B_d}(k')\biggr).
\]
For each fixed bin $B_d$,
\[
\sum_{k=0}^{p-1} \widehat{\mathbf{1}}_{B_d}(k)
= \sum_{k=0}^{p-1} \sum_{r \in B_d} e^{-2\pi i k r / p}
= \sum_{r \in B_d} \sum_{k=0}^{p-1} e^{-2\pi i k r / p}.
\]
Since $B_d \subset \{1, \ldots, p{-}1\}$, every $r \in B_d$
satisfies $r \not\equiv 0 \pmod p$. The standard
orthogonality of additive characters on
$\mathbb{Z}/p\mathbb{Z}$~\cite{iwaniec} gives
$\sum_{k=0}^{p-1} e^{-2\pi i k r / p} = 0$ for any such
$r$. Hence the inner sum vanishes for every $r \in B_d$, so
$\sum_k \widehat{\mathbf{1}}_{B_d}(k) = 0$ for each $d$.
The grand total is therefore a sum of products of zeros,
and equals zero.
\end{proof}

% ============================================================
\section{The Autocorrelation Formula}

\begin{theorem}\label{thm:autocorrelation}
For a prime $p$ where $b$ is a primitive root
($L = \operatorname{ord}_p(b) = p{-}1$), the cyclic
autocorrelation satisfies
\[
R(\ell) = \frac{1}{p}
\sum_{k=0}^{p-1}
G\!\left(k,\; -(k \cdot b^{-\ell}) \bmod p\right).
\]
\end{theorem}

\begin{proof}
By definition, the autocorrelation counts position-wise
digit matches between the repetend of $1/p$ and its rotation
by $\ell$, which (using $\delta(r) = \delta(s) \iff
\sum_d \mathbf{1}_{B_d}(r) \mathbf{1}_{B_d}(s) = 1$) is
\[
R(\ell) = \sum_{r=1}^{p-1} \sum_d
\mathbf{1}_{B_d}(r)\, \mathbf{1}_{B_d}(b^{\ell} r \bmod p).
\]
Apply the inverse additive Fourier expansion of each bin
indicator,
\[
\mathbf{1}_{B_d}(r) = \frac{1}{p} \sum_{k=0}^{p-1}
\widehat{\mathbf{1}}_{B_d}(k)\, e^{2\pi i k r / p},
\]
to both factors. After interchanging the order of summation
and collecting the $r$-dependence into a single phase,
\[
R(\ell) = \frac{1}{p^{2}}
\sum_{k, k'} G(k, k')
\sum_{r=1}^{p-1}
e^{2\pi i (k + k' b^{\ell}) r / p}.
\]

The inner sum is the standard nontrivial-character sum on
$(\mathbb{Z}/p\mathbb{Z})^*$, namely
\[
\sum_{r=1}^{p-1} e^{2\pi i m r / p}
=
\begin{cases}
p - 1 & \text{if } m \equiv 0 \pmod p, \\
-1 & \text{otherwise},
\end{cases}
\]
applied with $m = k + k' b^{\ell} \bmod p$. The condition
$k + k' b^{\ell} \equiv 0 \pmod p$ is equivalent (multiplying
by $b^{-\ell}$) to $k' \equiv -k b^{-\ell} \pmod p$. Define
\[
\Sigma_0(\ell) = \sum_{k = 0}^{p-1}
G\!\left(k,\, -k b^{-\ell} \bmod p\right),
\qquad
\Sigma_1(\ell) = \sum_{(k, k')} G(k, k')
\]
where the sum defining $\Sigma_1$ runs over the complementary
pairs $(k, k')$ with $k' \not\equiv -k b^{-\ell} \pmod p$.
Substituting,
\[
R(\ell) = \frac{1}{p^{2}}
\bigl[(p - 1)\, \Sigma_0(\ell) - \Sigma_1(\ell)\bigr].
\]

The two sums $\Sigma_0(\ell)$ and $\Sigma_1(\ell)$ partition
the full $(k, k')$-grid, so by Lemma~\ref{lem:G-sum-zero},
\[
\Sigma_0(\ell) + \Sigma_1(\ell)
= \sum_{k, k'} G(k, k') = 0,
\]
which gives $\Sigma_1(\ell) = -\Sigma_0(\ell)$.
Substituting back,
\[
R(\ell) = \frac{(p - 1)\, \Sigma_0(\ell)
+ \Sigma_0(\ell)}{p^{2}}
= \frac{p\, \Sigma_0(\ell)}{p^{2}}
= \frac{\Sigma_0(\ell)}{p}.
\qedhere
\]
\end{proof}

The cancellation $\Sigma_1(\ell) = -\Sigma_0(\ell)$ is the
key step. It follows from the vanishing total sum of $G$ and
eliminates the correction term entirely, leaving a clean
single sum. The numerical observation that $(p - 1) - (-1) = p$
exactly cancels the $p^2$ in the denominator is not coincidence
but a consequence of the dichotomy in the standard
nontrivial-character sum on $(\mathbb{Z}/p\mathbb{Z})^*$.

% ============================================================
\section{Worked Example: $p = 7$ in Base $10$}\label{sec:p7-example}

We verify Theorem~\ref{thm:autocorrelation} at $p = 7$ in
base $10$, where $b = 10$ has order $\operatorname{ord}_7(10) = 6
= p - 1$, so the primitive-root hypothesis holds. The bins of the
digit function are six singletons,
\[
B_1 = \{1\}, \;
B_2 = \{2\}, \;
B_4 = \{3\}, \;
B_5 = \{4\}, \;
B_7 = \{5\}, \;
B_8 = \{6\},
\]
and the bins $B_0, B_3, B_6, B_9$ are empty. Each nonempty
bin contains a single residue, so its Fourier coefficient
is the simple exponential
\[
\widehat{\mathbf{1}}_{B_d}(k) = e^{-2\pi i k r_d / 7},
\]
where $r_d$ is the residue of bin $d$. Summing over the six
nonempty bins,
\[
G(k, k') = \sum_{d \in \{1,2,4,5,7,8\}}
e^{-2\pi i (k + k') r_d / 7}
= \sum_{r=1}^{6} e^{-2\pi i (k + k') r / 7}.
\]
By the standard nontrivial-character sum,
\[
G(k, k') =
\begin{cases}
6 & \text{if } k + k' \equiv 0 \pmod 7, \\
-1 & \text{otherwise}.
\end{cases}
\]

Apply Theorem~\ref{thm:autocorrelation} with $\ell = 1$.
The inverse of $b = 10 \equiv 3 \pmod 7$ is $b^{-1} \equiv 5
\pmod 7$, so the orbit path is $k' = -5k \bmod 7$.
Tabulating $k + k' \bmod 7$ for $k = 0, 1, \ldots, 6$:
\begin{center}
\begin{tabular}{cccc}
\toprule
$k$ & $k' = -5k \bmod 7$ & $k + k' \bmod 7$ & $G(k, k')$ \\
\midrule
$0$ & $0$ & $0$ & $6$ \\
$1$ & $2$ & $3$ & $-1$ \\
$2$ & $4$ & $6$ & $-1$ \\
$3$ & $6$ & $2$ & $-1$ \\
$4$ & $1$ & $5$ & $-1$ \\
$5$ & $3$ & $1$ & $-1$ \\
$6$ & $5$ & $4$ & $-1$ \\
\bottomrule
\end{tabular}
\end{center}
The sum is $\Sigma_0(1) = 6 + 6 \cdot (-1) = 0$, so
$R(1) = 0/7 = 0$. Direct counting confirms this: the repetend
of $1/7$ in base $10$ is $142857$, and at lag $\ell = 1$ it
shifts to $428571$, which agrees with $142857$ at zero
positions. The same calculation at lags $\ell = 2, 3, 4, 5$
gives $\Sigma_0(\ell) = 0$ in each case, hence
$R(\ell) = 0$ for every $\ell \in \{1, 2, 3, 4, 5\}$, and
direct counting confirms this as well. Together with
$R(0) = 6$, this is the digit-partitioning autocorrelation
of $p = 7$ predicted by Corollary 2.2 of~\cite{paper7}.

% ============================================================
\section{The Complete Formula Chain}

The autocorrelation formula closes the chain from the digit
function to the eigenvalue spectrum:

\begin{enumerate}
\item \textbf{Bin sizes} $n_d$: determined by $p$ and $b$
      (the alignment limit paper~\cite{paper3}).
\item \textbf{Bin Fourier transforms}
      $\widehat{\mathbf{1}}_{B_d}(k)$: Dirichlet kernels with
      phases from bin starting positions (the spectral power
      paper~\cite{paper8}).
\item \textbf{Cross-spectral function} $G(k,k')$: products of
      bin transforms (the present paper).
\item \textbf{Autocorrelation} $R(\ell) = \Sigma_0(\ell)/p$:
      $G$ evaluated along the orbit path
      $k' = -kb^{-\ell} \bmod p$ (the present paper).
\item \textbf{Eigenvalues}
      $\lambda_j = (1/L) \sum_{\ell} R(\ell)\, e^{2\pi i
      j\ell/L}$: DFT of $R$ (the spectral structure
      paper~\cite{paper7}).
\end{enumerate}

Every step is explicit and computable from $p$ and $b$ alone.
The bin starting positions, which appear as the phase factors
$e^{-2\pi i k a_d / p}$ in $\widehat{\mathbf{1}}_{B_d}(k)$, are
determined by the bin sizes through cumulative summation
($a_d = 1 + \sum_{d' < d} n_{d'}$), so the entire chain is a
function of the pair $(p, b)$ alone. No empirical fitting, no
orbit enumeration, no position-by-position digit comparison is
required.

% ============================================================
\section{The Multi-Orbit Case}\label{sec:multi-orbit}

When $L < p - 1$, the orbit $\langle b \rangle$ is a proper
subgroup of $(\mathbb{Z}/p\mathbb{Z})^{*}$, and $R(\ell)$ sums
only over the $L$ orbit elements. The primitive-root proof
does not extend directly: the sum over the restricted orbit
involves additive characters evaluated on a multiplicative
subgroup (incomplete Gauss sums), which do not reduce to the
two-valued dichotomy used above.

The correct generalization requires a subgroup projector
involving multiplicative characters of the quotient group
$(\mathbb{Z}/p\mathbb{Z})^{*} / \langle b \rangle$.
This is not pursued here. The multi-coset autocorrelation
formula remains open.

% ============================================================
\section{Remarks}

\subsection*{G as the complete invariant}

For primitive-root primes, the cross-spectral function
$G(k, k')$ determines the autocorrelation $R(\ell)$ and hence
the eigenvalue spectrum of the cross-alignment matrix.
The function $G$ is symmetric in its arguments,
$G(k, k') = G(k', k)$, since the defining product
$\widehat{\mathbf{1}}_{B_d}(k) \widehat{\mathbf{1}}_{B_d}(k')$
is commutative. It is not Hermitian in general, however:
$\overline{G(k', k)} = G(-k', -k) \ne G(k, k')$ unless the
bin partition is symmetric under the involution
$r \mapsto -r$ on $(\mathbb{Z}/p\mathbb{Z})^*$. The
spectral power $\Phi(k) = G(k, -k)$ from~\cite{paper8}
captures the magnitude of the bin Fourier coefficients.
The off-diagonal entries $G(k, k')$ for $k \ne -k$ capture
the phase relationships. Together, the on- and off-diagonal
data determine $R(\ell)$ via Theorem~\ref{thm:autocorrelation}
and through it the eigenvalues of $\mathbf{A}(p)$.

\subsection*{The orbit path in $(k,k')$ space}

The map $\ell \mapsto (k, -kb^{-\ell} \bmod p)$ traces a path
through the $(k,k')$ plane for each fixed $k$. The
autocorrelation $R(\ell)$ is the sum of $G$ along these paths.
Different primes produce different path geometries, which is why
the eigenvalue spectrum varies across primes even when the bin
sizes are similar.

\subsection*{From character sums to the program}

The formula $R(\ell) = \Sigma_0(\ell)/p$ is a character sum: it
evaluates the cross-spectral function at points determined by the
multiplicative structure of $(\mathbb{Z}/p\mathbb{Z})^{*}$.
This confirms the connection between the fractional field's
coherence geometry and harmonic analysis suggested in earlier
papers, and completes the bridge from the digit function's
Euclidean structure to the additive character theory of cyclic
groups~\cite{terras}.  The classical theory of character sums and
their applications to arithmetic functions is developed
in~\cite{hardy,montgomery}.

% ============================================================
\begin{thebibliography}{9}

\bibitem{paper3}
A.~S.~Petty,
The alignment limit for all primes,
research note, 2020.

\bibitem{paper6}
A.~S.~Petty,
The cross-alignment matrix,
research note, 2021.

\bibitem{paper7}
A.~S.~Petty,
The spectral structure of fractional fields,
research note, 2021.

\bibitem{paper8}
A.~S.~Petty,
The spectral power of the digit function,
research note, 2021.

\bibitem{iwaniec}
H.~Iwaniec and E.~Kowalski,
\emph{Analytic Number Theory},
Amer.\ Math.\ Soc., 2004.

\bibitem{terras}
A.~Terras,
\emph{Fourier Analysis on Finite Groups and Applications},
Cambridge University Press, 1999.

\bibitem{hardy}
G.~H.~Hardy and E.~M.~Wright,
\emph{An Introduction to the Theory of Numbers},
6th~ed., Oxford University Press, 2008.

\bibitem{montgomery}
H.~L.~Montgomery and R.~C.~Vaughan,
\emph{Multiplicative Number Theory I: Classical Theory},
Cambridge University Press, 2007.

\end{thebibliography}

\end{document}
